\section{Anatomie du repo implicite}\label{sec-anatomy}

\subsection{L'écart de panel}
Pour un pilier $T$, on note $\bA(T)$ et $\DA(T)$ les moyennes des vues de repo et de dividendes sur le panel options,
et $\DB(T)$ le consensus dividendes du panel dividendes après allocation. Les vues sont des perturbations d'un centre
commun, $b_j=b_0+\delta b_j$ et $D_j=D_0+\delta D_j$, d'ordre $\delta$.

\begin{proposition}[Décomposition par écart de panel]\label{prop-gap}
Sous les \cref{ass-consistent,ass-panels,ass-consensus},
\begin{equation}\label{eq-gap}
b^\T(T)=\bA(T)+\frac{\DA(T)-\DB(T)}{T\,\bigl(S-D_0(T)\bigr)}+O(\delta^2).
\end{equation}
\end{proposition}
\begin{proof}
Pour $j\in\Aset$, le \cref{lem-additivity} donne $\ln F_j=\ln F_0-T\,\delta b_j-\delta D_j/(S-D_0)+O(\delta^2)$. Le
forward de consensus est la moyenne des $F_j$, et le logarithme d'une moyenne égale la moyenne des logarithmes à un
terme de variance d'ordre $\delta^2$ près, d'où
$\ln\overline{F}=\ln F_0-T\,\overline{\delta b}_{\Aset}-\overline{\delta D}_{\Aset}/(S-D_0)+O(\delta^2)$. Le contrôleur
alloue le consensus dividendes, $\DB=D_0+\overline{\delta D}_{\Bset}$, et applique le bootstrap :
$B^\T=R-\ln\overline{F}+\ln(S-\DB)
=B_0+T\,\overline{\delta b}_{\Aset}+\bigl(\overline{\delta D}_{\Aset}-\overline{\delta D}_{\Bset}\bigr)/(S-D_0)+O(\delta^2)$.
On divise par $T$ et on utilise $b_0+\overline{\delta b}_{\Aset}=\bA$.
\end{proof}

Le repo implicite est le repo propre au panel options plus un \emph{écart de panel} : la différence entre la vue
dividendes des contributeurs options, jamais observée, et celle du panel dividendes, divisée par $T(S-D)$. Pour un
contributeur donné, les erreurs sur le forward et sur les dividendes se compensent dans le repo implicite ; entre deux
panels, elles ne se compensent pas.

\begin{corollary}[Taille de l'écart]\label{cor-gap}
Supposons les erreurs de vue dividendes en $T$ i.i.d.\ de variance $\sigma_D^2(T)$ dans les deux populations, et
supposons que les deux panels partagent $m$ contributeurs qui utilisent la même vue dividendes dans les deux services.
Alors
\begin{equation}\label{eq-gap-sd}
\sd\bigl[b^\T-\bA\bigr]=\frac{\sigma_D(T)}{T\,(S-D(T))}\sqrt{\frac1{n_A}+\frac1{n_B}-\frac{2m}{n_An_B}}.
\end{equation}
\end{corollary}
\begin{proof}
$\DA-\DB=\sum_j(a_j-c_j)\,\delta D_j$ avec $a_j=1/n_A$ sur $\Aset$, $c_j=1/n_B$ sur $\Bset$ et zéro ailleurs. Sa variance
vaut $\sigma_D^2\sum_j(a_j-c_j)^2=\sigma_D^2\bigl[m(1/n_A-1/n_B)^2+(n_A-m)/n_A^2+(n_B-m)/n_B^2\bigr]$, qui se simplifie
en $\sigma_D^2\,(1/n_A+1/n_B-2m/(n_An_B))$.
\end{proof}

Avec des panels identiques ($m=n_A=n_B$), l'écart disparaît ; avec des panels disjoints, il est maximal. La
dispersion $\sigma_D(T)$ d'un dividende cumulé croît avec $T$, à peu près comme $T$ ou plus vite : l'écart est donc
faible en portage sur le court terme mais ne diminue pas en repo sur le long terme.

\begin{example}[Taille de l'écart pour une action]\label{ex-gap}
Spot $100$, rendement en dividendes de $4\,\%$, dispersion des contributeurs de $1\,\%+3\,\%\cdot T$ du dividende
cumulé, tailles de panel de la catégorie synthétique « action UE ». La formule \eqref{eq-gap-sd} avec $m=0$ et une
simulation contributeur par contributeur sans désaccord sur le repo donnent
\begin{center}\input{generated/\lang/tab_gap}\end{center}
À un an, si six contributeurs étaient communs aux deux services ($m=6$, $n_A=15$, $n_B=12$), l'écart passerait de
$6{,}5$ à $4{,}8$\,pb.
\end{example}

\begin{figure}[ht]
\centering\includegraphics[width=\linewidth]{figures/\lang/panel_gap.pdf}
\caption{(a) Bruit de repo dû au seul écart de panel, \cref{eq-gap-sd}, pour trois configurations de panels. (b) Un
décalage de spot $\varepsilon$ est constant en portage et se lit donc $\varepsilon/T$ en repo (\cref{prop-spot}).}
\label{fig-gap}
\end{figure}

\subsection{Sortie de contributeurs}
\begin{proposition}[Moyennes de panels emboîtés]\label{prop-nested}
Soit $e_{jk}$ l'erreur du contributeur $j$ à la maturité $k$, indépendante d'un contributeur à l'autre, avec
$\Var e_{jk}=\sigma_k^2$ et $\Cov(e_{jk},e_{jl})=\rho_{kl}\sigma_k\sigma_l$. Si les panels sont emboîtés, de tailles
$n_k$, les moyennes $\bar e_k$ vérifient
\begin{equation}\label{eq-nested}
\Cov(\bar e_k,\bar e_l)=\frac{\rho_{kl}\,\sigma_k\sigma_l}{\max(n_k,n_l)}.
\end{equation}
En particulier, avec des vues persistantes ($\rho=1$, $\sigma_k=\sigma_l=\sigma$) et $n_l\le n_k$,
$\sd(\bar e_l-\bar e_k)=\sigma\sqrt{1/n_l-1/n_k}$.
\end{proposition}
\begin{proof}
$\Cov(\bar e_k,\bar e_l)=\frac{1}{n_kn_l}\sum_{j\le\min(n_k,n_l)}\rho_{kl}\sigma_k\sigma_l
=\frac{\min(n_k,n_l)}{n_kn_l}\rho_{kl}\sigma_k\sigma_l$, et $\min/(n_kn_l)=1/\max$. Ensuite
$\Var(\bar e_l-\bar e_k)=\sigma^2(1/n_l+1/n_k-2/n_k)=\sigma^2(1/n_l-1/n_k)$.
\end{proof}

\begin{example}
En unités de dispersion des contributeurs, le consensus saute de lui-même de $\vJumpA\,\sigma$ quand un panel passe
de $16$ à $12$ contributeurs, de $\vJumpB\,\sigma$ de $12$ à $7$, de $\vJumpC\,\sigma$ de $15$ à $7$ et de
$\vJumpD\,\sigma$ de $10$ à $4$. Un saut de cette taille lors d'une sortie de contributeurs est du bruit, pas une
information ; le modèle de bruit de la \cref{sec-noise} le reproduit grâce à \cref{eq-nested}.
\end{example}

\subsection{Désaccords sur la date de détachement}
\begin{proposition}[Désaccord de date]\label{prop-timing}
Soit un dividende $d_i$ détaché en $t_i\in(T_k,T_{k+1})$ dans l'allocation des dividendes, et supposons qu'une fraction
$p$ du panel options le place avant $T_k$ (tous le placent avant $T_{k+1}$). En notant $\tilde d_i=d_i/G(t_i)$, au premier
ordre :
\begin{enumerate}[label=(\roman*)]
\item l'écart de portage ne bouge qu'en $T_k$ : $\delta y_k=p\,\tilde d_i/(S-D_k)$, soit un pic de repo
$p\,\tilde d_i/\bigl(T_k(S-D_k)\bigr)$ ;
\item en repo forward entre piliers, c'est un dipôle : $+\delta y_k/(T_k-T_{k-1})$ puis $-\delta y_k/(T_{k+1}-T_k)$ ;
\item la variance des log-forwards des contributeurs en $T_k$ augmente de $p(1-p)\,\tilde d_i^{\,2}/(S-D_k)^2$.
\end{enumerate}
\end{proposition}
\begin{proof}
(i) En $T_k$, une fraction $p$ du panel options inclut le dividende, donc $\ln\overline{F}_k$ baisse de
$p\,\tilde d_i/(S-D_k)$ (\cref{lem-additivity}) tandis que le $D^\T_k$ alloué l'exclut ; par \cref{eq-additivity},
$B^\T_k$ augmente d'autant. Avant $T_k$ personne n'inclut le dividende, après $T_{k+1}$ tout le monde l'inclut :
aucun autre pilier ne bouge. (ii) Le repo forward est la pente de $B$ entre piliers. (iii) Formule de la variance
totale : les contributeurs forment deux groupes dont les log-forwards moyens diffèrent de $\tilde d_i/(S-D_k)$, de poids
$p$ et $1-p$.
\end{proof}

\begin{example}\label{ex-timing}
Dans l'exemple de la \cref{fig-gallery}(c), un dividende de $2{,}18$ est détaché $17$ jours après l'échéance à six mois,
et $37\,\%$ du panel options le placent avant. Le pic vaut $0{,}37\times2{,}18/(0{,}5\times98)\approx1{,}6\,\%$, soit
environ $165$\,pb de repo sur un seul pilier. La dispersion des contributeurs à cette échéance est presque dix fois
celle de ses voisines, conformément à (iii).
\end{example}

\subsection{Profil de détachement périmé}
\begin{proposition}[Décalage de profil]\label{prop-profile}
Supposons que les contributeurs options utilisent les vraies dates $t_i$, tandis que le profil d'allocation les
place en $t_i+\Delta$. Alors, pour $\Delta>0$ (symétriquement pour $\Delta<0$),
\[
\delta y_k=\sum_{i:\,T_k\in[t_i,\,t_i+\Delta)}\frac{\tilde d_i}{S-D_k}\,:
\]
seuls bougent les piliers situés entre une vraie date de détachement et sa copie décalée.
\end{proposition}
\begin{proof}
$D^\T_k-D_k=\sum_i\tilde d_i\bigl(\one\{t_i+\Delta\le T_k\}-\one\{t_i\le T_k\}\bigr)$, qui vaut $-\tilde d_i$ exactement
lorsque $t_i\le T_k<t_i+\Delta$ ; on applique \cref{eq-additivity}.
\end{proof}
Pour les payeurs saisonniers, les piliers touchés reviennent chaque année (par exemple les échéances d'avril des
actions européennes), ce qui produit le zigzag que les contrôleurs reconnaissent. Les fins de période ne sont jamais
touchées, car les totaux par période sont justes.

\subsection{Erreurs sur une période de dividendes}
\begin{proposition}[Erreur de période]\label{prop-bucket}
Si le niveau Totem de la période $y$ est erroné de $\delta L$, alors $\delta y_k=-A_{ky}\,\delta L/(S-D_k)$. Pour une
période à une seule date de détachement, c'est une marche qui commence au premier pilier après le détachement ; en
repo, c'est un saut qui décroît comme $1/T$.
\end{proposition}
\begin{proof}
Le \cref{lem-allocation} donne $\delta D_k=A_{ky}\delta L$ ; on applique \cref{eq-additivity} à forwards inchangés.
$A_{ky}$ est nul avant le détachement et constant ($\approx1/G(t_i)$) après.
\end{proof}

\subsection{Décalage du spot et dividendes proches}
\begin{proposition}[Décalage du spot]\label{prop-spot}
Si les contributeurs options valorisent avec un spot $S(1+\varepsilon)$ alors que le contrôleur utilise $S$,
\begin{equation}\label{eq-spot}
y_k=-\ln\Bigl(1+\frac{\varepsilon S}{S-D_k}\Bigr)\approx-\frac{\varepsilon S}{S-D_k}\qquad\text{pour tout }k,
\end{equation}
quasi constant d'un pilier à l'autre (exactement constant sans dividende), donc une erreur de repo proche de
$-\varepsilon/T_k$. Un dividende antérieur au premier pilier, erroné de $\delta d$, produit la même signature,
$-\delta\tilde d/(S-D_k)$.
\end{proposition}
\begin{proof}
Le forward de chaque contributeur devient $G(S(1+\varepsilon)-D)$, donc $\ln\overline{F}_k$ varie de
$\ln\bigl((S(1+\varepsilon)-D_k)/(S-D_k)\bigr)$ ; par \cref{eq-additivity}, $B^\T_k$ varie de l'opposé. Pour le
dividende proche, $D_k$ varie de $\delta\tilde d$ à tous les piliers.
\end{proof}

\subsection{Définitions brut/net et biais proportionnels}
\begin{proposition}[Biais proportionnel sur les dividendes]\label{prop-gross}
Si chaque dividende Totem vaut $(1+\tau)$ fois le dividende contenu dans les forwards (par exemple brut contre net de
retenue à la source), alors $y_k\approx-\tau D_k/(S-D_k)$. Pour un rendement continu $q$, $D_k\approx SqT_k$ et
$y_k\approx-\tau q\,T_k$, qui a exactement la forme d'un écart de niveau de repo de $-\tau q$.
\end{proposition}
\begin{proof}
$\delta D_k=\tau D_k$ dans \cref{eq-additivity} ; pour un rendement, $D_k=S(1-e^{-qT_k})\approx SqT_k$ et $S-D_k\approx S$.
\end{proof}

\subsection{Ce que les données ne peuvent pas dire}
\begin{theorem}[Non-identifiabilité de la répartition]\label{thm-nonid}
Supposons que les prix d'options européennes ne dépendent des vues d'un contributeur qu'à travers ses forwards. Pour
toute transformation au niveau des contributeurs $(b_j,D_j)\mapsto(b_j+\delta_j,D'_j)$ avec
$\delta_j(T)\,T=\ln(S-D'_j(T))-\ln(S-D_j(T))$, tous les prix d'options, et donc toutes les statistiques du panel options
(consensus, dispersion, nombre), sont inchangés ; les statistiques du panel dividendes le sont aussi. Le repo implicite
$b^\T$, fonction de ces statistiques, est donc inchangé, alors que le repo du panel options $\bA$ varie de la moyenne
des $\delta_j$. Par conséquent, l'écart de panel $b^\T-\bA$ n'est pas une fonction des données observables.
\end{theorem}
\begin{proof}
La transformation laisse $\ln F_j=R-B_j+\ln(S-D_j)$ inchangé pour tout $j$ et tout $T$ ; les prix d'options de chaque
contributeur sont donc inchangés, ainsi que toute statistique calculée à partir d'eux. Le panel dividendes n'est pas
touché. $b^\T$ n'est calculé qu'à partir de ces statistiques, il est donc inchangé, alors que $\bA$ se déplace de la
moyenne des $\delta_j$, qui peut prendre n'importe quelle valeur.
\end{proof}

\begin{corollary}
Aucun estimateur construit à partir des seules données de consensus Totem (lissage de courbe, regroupement, réseaux de
neurones) ne peut supprimer l'écart de panel. Il reste possible d'en modéliser la taille (\cref{sec-noise}), de
reconnaître les mécanismes structurés à leur signature (\cref{sec-detection}) et de se passer de la répartition là où
l'exposition ne l'exige pas (\cref{sec-control}).
\end{corollary}

\begin{remark}[Ce qui lève l'indétermination]
Les options américaines portent un peu d'information sur la répartition à travers l'exercice anticipé ; des
contributeurs communs réduisent l'écart (\cref{cor-gap}) ; des ancres externes fixent directement l'un des deux
côtés : contrats à terme cotés sur dividendes, contrats à terme de rendement total qui cotent un spread de financement,
commissions de prêt-emprunt de titres pour le court terme des actions.
\end{remark}

\begin{insight}
Résumé des signatures dans l'espace de portage (gabarits du détecteur entre crochets) :
\begin{center}\small
\begin{tabular}{p{4.3cm}p{4.9cm}p{4.4cm}}
\toprule
Mécanisme & Écart de portage $y$ & Écart de repo\\
\midrule
Écart de panel (\ref{prop-gap}) & bruit lisse plus incréments aux détachements & bruit à peu près plat en maturité\\
Sortie de contributeurs (\ref{prop-nested}) & marche à la sortie [marche] & saut décroissant en $1/T$\\
Désaccord de date (\ref{prop-timing}) & un pilier [pic] & pic, dipôle en repo forward\\
Profil périmé (\ref{prop-profile}) & piliers en saison [saison, pic] & zigzag\\
Erreur de période (\ref{prop-bucket}) & marche après le détachement [marche] & saut décroissant en $1/T$\\
Décalage du spot (\ref{prop-spot}) & constante [niveau] & hyperbole $\varepsilon/T$\\
Biais proportionnel (\ref{prop-gross}) & rampe de marches ; $\approx\beta T$ pour un rendement & marches ; un niveau pour un indice\\
Erreur de marquage FO & $\delta b(T)\,T$ lisse [modèle nul] & niveau ou pente lisse\\
\bottomrule
\end{tabular}
\end{center}
\end{insight}
